\section*{群の直積}
$G_1, G_2$を群とする. これらは直積$G_1\times G_2$の部分群とみなせる.
\begin{tcolorbox}
	\begin{prp*}
		(1) $G_1\times G_2$の中で$G_1$の元と$G_2$の元は可換である.

		(2) $G_1, G_2$は$G_1\times G_2$の正規部分群である.
	\end{prp*}
\end{tcolorbox}

\begin{proof}
	(1) $(g_1, 1_{G_2})(1_{G_1},g_2)=(1_{G_1},g_2)(g_1,1_{G_2})=(g_1,g_2)$
	となるので,$G_1$の元と$G_2$の元は可換である.

	(2) $g_1 \in G_1, (g'_1, g'_2)\in G_1\times G_2$なら,
	\[ (g'_1, g'_2)(g_1,1_{G_2})(g'_1,g'_2)^{-1}
	=(g'_1 g_1 {g'_1}^{-1},1_{G_2})\]
	なので,$G_1 \lhd G_1 \times G_2$である.
	同様に$G_2 \lhd G_1 \times G_2$である.
\end{proof}

\begin{tcolorbox}
	\begin{prp*}
		$G$が群,$H,K\in G$が正規部分群で$H\cap K=\{1_G\}, HK=G$とする.
		このとき, $G$は直積$H\times K$と同型である.
	\end{prp*}
\end{tcolorbox}
\begin{proof}
  写像$\phi:H\times K\to G$を$\phi(h,k)=hk$と定義する.
  $HK=G$なので,任意の$g\in G$に対して, $g=hk$となるような$h\in H, k\in K$が存在する.
  したがって$\phi$は全射である.
  
  $H,K$は$G$の正規部分群なので,$khk^{-1}\in H$であり, また$h^{-1}\in H$であるから,
  \[(khk^{-1})h^{-1}=khk^{-1}h^{-1}\in H\]
  となる.
  同様に, $hk^{-1}h^{-1}\in K$で
  \[k(hk^{-1}h^{-1})=khk^{-1}h^{-1}\in K\]
  である.
  $H\cap K=\{1_G\}$なので, $khk^{-1}h^{-1}=1_G$.よって$kh=hk$であり,
  $h'\in H, k'\in K$とすると
  \[ \phi(h,k)\phi(h',k')=hkh'k'=hh'kk'=\phi(hh',kk').\]
  したがって$\phi:H\times K\to G$は準同型である.

  $\phi(h,k)=hk=1_G$とすると,$h=k^{-1}$であり,$h\in H$かつ$h\in K$なので
  $h=k^{-1}\in H\cap K$である.$H\cap K=\{1_G\}$であるから,
  $\mathrm{Ker}\phi = \{1_G\}$.したがって$\phi: H\times K\to G$は単射である.

  以上より, $G\cong H\times K$である.
\end{proof}

\begin{tcolorbox}
	\begin{thm*}
		（中国式剰余定理）$m,n \neq 0$が互いに素な整数なら,
		$\mathbb{Z}/mn\mathbb{Z}$は
		$\mathbb{Z}/m\mathbb{Z}\times \mathbb{Z}/n\mathbb{Z}$
		と同型である.
	\end{thm*}
\end{tcolorbox}

\begin{proof}
	$\mathbb{Z}/mn\mathbb{Z}$から
	$\mathbb{Z}/m\mathbb{Z}\times \mathbb{Z}/n\mathbb{Z}$
	への写像$\phi$を
	\[ \phi(x+mn\mathbb{Z})=(x+m\mathbb{Z},x+n\mathbb{Z})\]
	と定義する.
	$y\in x+mn\mathbb{Z}$なら,
	$a\in \mathbb{Z}$により$y=x+mna$となる.
	よって,
	$y\in x+m\mathbb{Z},x+n\mathbb{Z}$であり,
	$y+m\mathbb{Z}=x+m\mathbb{Z}, y+n\mathbb{Z}=x+n\mathbb{Z}$となる.
	したがって, 写像$\phi$はwell-definedである.
\end{proof}
